\section{Conclusion}

This paper has established a precise and strictly internal notion of
falsifiability for the enterprise-as-continuum model formulated within the
Ontology of Continua (OC Core) and its executable specification
ETS.K\_EA.

\medskip

Falsifiability is shown to be neither empirical nor operational in the
classical sense.
It is defined exclusively as an internal logical property: an admissible
observation falsifies the model if and only if it renders the set of
compatible model realizations \(\Omega(K_{EA})\) empty.
Equivalently, falsification occurs precisely when the constraints induced
by an admissible observation are internally inconsistent with the mandatory
axioms, invariants, thresholds, or structural constraints that define the
enterprise continuum.
This criterion provides a necessary and sufficient characterization of
falsification at the level of formal existence.

\medskip

The analysis explicitly separates falsification from several neighboring
phenomena that are frequently conflated with it.
Non-identifiability, underdetermination, and the coexistence of multiple
compatible model realizations are admissible outcomes and do not constitute
falsification.
Likewise, logical undecidability reflects limits of inference, computation,
or expressibility rather than inconsistency of the model.
The \emph{STOP} outcome is formally characterized as a first-class
diagnostic termination condition marking the exhaustion of admissible
inference paths under constraint preservation, not a failure or refutation
of the model.

\medskip

Admissibility is shown to be logically prior to falsifiability.
Only observations that are expressible within the formal vocabulary and
constraints of ETS.K\_EA and that do not presuppose violations of mandatory
model constraints can participate in a falsifiability assessment.
Inputs outside this domain are excluded before any question of
falsification arises and are therefore neither confirming nor disconfirming
by construction.

\medskip

Taken together, these results place the enterprise-as-continuum model in a
narrow and well-defined logical position.
The model is falsifiable in principle, but only through demonstrable
internal inconsistency within its formal state space.
All other apparent failure modes—ambiguity, undecidability, diagnostic
termination, or external disagreement—are explicitly classified as
non-falsifying outcomes of the formalism.

\medskip

This closure completes the formal characterization of what it means for the
enterprise-as-continuum model to be falsifiable, what it explicitly does not
claim, and why these boundaries are necessary for logical coherence.
No interpretive, empirical, or prescriptive conclusions are implied beyond
the formal scope established here.
